\SourceLeaf{92}tus magneticus etiam absque perfecto aequilibrio inter illas vires durabilis esse potest, si modo vires fortiores extraneae inde arceantur. Etiamsi caussa huius phaenomeni ignota sit, tamen eam ita imaginari licet, ac si partes ponderabiles corporis secundi generis motui fluidorum magneticorum cum ipsis iunctorum aliquod obstaculum frictioni simile opponant, quae resistentia in ferro molli vel nulla est, vel saltem perparva.

In disquisitione theoretica hi duo casus tractationem prorsus diversam requirunt, sed in commentatione praesente de solis corporibus secundi generis sermo erit: in experimentis, de quibus agemus, stabilitas status magnetici in singulis corporibus ad illa adhibitis erit suppositio fundamentalis, probeque proin cavendum est, ne inter experimenta alia corpora, quae hunc statum mutare possent, nimis prope accedant.

Attamen exstat quaedam caussa mutationis, cui etiam corpora secundi generis obnoxia sunt, puta calor. Nimirum experientia docet, statum magneticum corporis variari cum eius temperatura, caloremque auctum intensitatem magnetismi debilitare, ita tamen, ut nisi corpus ultra modum calefactum fuerit, cum priori temperatura prior quoque status magneticus redeat. Haec dependentia per experimenta idonea determinanda est, et si operationes ad idem experimentum pertinentes sub temperaturis inaequalibus institutae sunt, ante omnia ad eandem temperaturam revocandae erunt.

\Article{4}
Independenter a viribus magneticis, quas corpora singularia satis sibi vicina in se mutuo exercere videmus, alia vis in fluida magnetica agit, quam quum ubique terrarum se manifestet, ipsi globo terrestri tribuimus, atque magnetismum terrestrem vocamus. Duplici modo haec vis se exserit: corpora secundi generis, in quibus magnetismus evolutus est, si in centro gravitatis sustinentur, ad directionem determinatam sollicitantur: contra in corporibus primi generis fluida magnetica per istam vim sponte separantur, quae separatio, si corpora figurae idoneae eliguntur atque in positione idonea collocantur, persensibilis reddi potest. Utrumque phaenomenon explicatur, vim illam ita concipiendo, ut fluidum magneticum boreale in quovis loco versus certam directionem propellat, australe vero aequali intensitate versus oppositam. Directio prior semper intelligitur, dum de directione magnetismi terrestris loquimur, quae proin per inclinationem ad pla\SourceLeaf{93}num horizontale atque declinationem plani verticalis, in quo agit, a plano meridiano determinatur: illud \textit{planum meridianum magneticum} vocatur. Intensitas autem magnetismi terrestris per vim motricem, quam in unitatem fluidi magnetici liberi exserit, aestimanda est.

Haec vis non modo in diversis terrae locis diversa est, sed etiam in eodem loco variabilis, tum per saecula et annos. tum per anni aestates dieique horas. Respectu directionis haec variabilitas dudum quidem nota fuit: sed respectu intensitatis hactenus tantummodo per horas diei animadverti potuit, quum subsidiis ad longiora temporis intervalla aptis caruissemus. Huic defectui in posterum reductio intensitatis ad mensuram absolutam remedium afferet.

\Article{5}
Ut actio magnetismi terrestris in corpora magnetica secundi generis (qualia semper abhinc subintelligenda sunt) calculo subiiciatur, concipiatur tale corpus in partes infinite parvas divisum, sitque $\mathrm{d}m$ elementum magnetismi liberi in particula, cuius coordinatae respectu trium planorum inter se normalium et respectu corporis fixorum denotentur per $x,y,z$: elementa fluidi australis negative accipi supponimus. Ita primo patet, integrale $\int\mathrm{d}m$ per totum corpus collectum (imo per quamlibet corporis partem mensurabilem) esse $=0$. Statuamus $\int x\mathrm{d}m=X$, $\int y\mathrm{d}m=Y$, $\int z\mathrm{d}m=Z$, quae quantitates vocari poterunt momenta magnetismi liberi respectu trium planorum fundamentalium, sive respectu axium in ipsa normalium. Quum denotante $a$ quantitatem constantem arbitrariam, fiat $\int(x-a)\mathrm{d}m=X$, patet, momentum respectu axis dati pendere tantummodo ab eius directione, non autem ab eius initio. Si per initium coordinatarum axem quartum ducimus, qui cum primariis faciat angulos $A,B,C$, momentum elementi $\mathrm{d}m$ respectu huius axis erit $=(x\cos A+y\cos B+z\cos C)\mathrm{d}m$, adeoque momentum magnetismi liberi in toto corpore
\[
=X\cos A+Y\cos B+Z\cos C=V
\]
Statuatur
\[
\BareRoot{(XX+YY+ZZ)}=M,\enspace\text{atque }X=M\cos\alpha,\enspace Y=M\cos\beta,\quad Z=M\cos\gamma
\]
ducaturque axis quintus, qui cum tribus primariis faciat angulos $\alpha,\beta,\gamma$, et cum axi quarto angulum $\omega$; unde quum constet esse
\SourceLeaf{94}
\[
\cos\omega=\cos A\cos\alpha+\cos B\cos\beta+\cos C\cos\gamma,\quad\text{fiet }V=M\cos\omega
\]
Hunc axem quintum simpliciter vocamus corporis \textit{axem magneticum}, eiusque \textit{directionem} ad valorem positivum radicalis $\BareRoot{(XX+YY+ZZ)}$ referri supponimus. Si axis quartus cum hoc axe magnetico coincidit, momentum $V$ fit $=M$, quod manifesto inter omnia est maximum: momentum respectu cuiuslibet alius axis invenitur, multiplicando hoc momentum maximum (quod quoties ambiguitas non metuenda est, simpliciter momentum magnetismi vocari potest) per cosinum anguli inter hunc axem atque axem magneticum. Momentum respectu cuiusvis axis in axem magneticum normalis fit $=0$, negativum vero respectu cuiusvis axis, qui cum axe magnetico angulum obtusum facit.

Axis itaque magneticus non est recta determinata, quum per punctum quodlibet duci possit, sed tantummodo directio determinata, sive adsunt infinite multi axes magnetici inter se paralleli. E quibus si aliquem ad lubitum eligimus, longitudinemque determinatam ipsi tribuimus, eius termini vocantur poli, alter australis, a quo, alter borealis, versus quem directio axis procedit.

\Article{6}
Si in singulas fluidorum magneticorum particulas vis agit intensitate et directione constans, vis totalis in corpus inde resultans facile e principiis staticis derivatur, quum in corporibus, quae hic consideramus, particulae illae fluiditatem quasi amiserint, et cum corpore ponderabili massam unam rigidam sistant. Agat in quamvis moleculam magneticam $\mathrm{d}m$ vis motrix $=P\mathrm{d}m$ secundum directionem $D$ (ubi pro moleculis fluidi australis signum negativum iam per se directionem oppositam implicat); sint $A,B$ duo corporis puncta in directione axis magnetici iacentia, eorumque distantia $=r$, positive accepta, dum axis magneticus tendit ab $A$ versus $B$: ita facile intelligitur, si viribus istis duae novae adiungantur, utraque $=\frac{PM}{r}$, et quarum altera agat in $A$ secundum directionem $D$, altera in $B$ secundum directionem oppositam, inter omnes has vires aequilibrium fore. Quapropter vires priores aequivalebunt duabus viribus $=\frac{PM}{r}$, quarum altera in $B$ secundum directionem $D$, altera in $A$ secundum directionem oppositam agit, manifestoque hae duae vires in unam conflari nequeunt.

Si praeter vim $P$ alia similis $P'$ secundum directionem $D'$ in corporis fluida magnetica agit, eius loco iterum duae aliae vel in eadem puncta $A,B$, vel ge\SourceLeaf{95}neralius in puncta alia $A',B'$ agentia substitui possunt, dummodo $A'B'$ quoque sit axis magneticus, et quidem faciendo distantiam $A'B'=r'$, hae vires debent esse $=\frac{P'M}{r'}$, atque in $B'$ agetur secundum directionem $D'$, in $A'$ secundum oppositam, et perinde de pluribus.

Vi magneticae terrestri intra tam parvum spatium, quantum corpus experimentis subiiciendum explet, tuto intensitatem atque directionem ubique constantem (etiamsi respectu temporis variabilem) tribuere, adeoque ea, quae modo diximus, ad eam applicare licet. Sed commodum esse potest, statim ab initio in duas vires eam resolvere, alteram horizontalem $=T$, alteram verticalem, nostris regionibus deorsum tendentem, $=T'$. Quum, si pro posteriori duas alias in puncta $A',B'$ agentia substituere placet, tum punctum $A'$ tum distantiam $A'B'=r'$ pro lubitu assumere liceat, pro $A'$ adoptabimus centrum gravitatis, et denotato pondere corporis, i. e. vi motrice, quam gravitas massae corporis inducit, per $p$, statuemus $\frac{T'M}{p}=r'$. Hoc pacto effectus vis $T'$ resolvitur in vim $=p$ in $A'$ sursum, atque in aliam aequalem in $B'$ deorsum tendentem, adeoque quum prior manifesto per ipsam gravitatem destruatur, effectus vis magneticae terrestris verticalis simpliciter reducitur ad transpositionem centri gravitatis ex $A'$ in $B'$. Ceterum manifestum est, pro iis regionibus, ubi vis magnetica terrestris facit angulum acutum cum linea verticali, sive ubi eius pars verticalis fluidum magneticum boreale sursum propellit, similem transpositionem centri gravitatis in axi magnetico versus polum australem locum habere.

Ex hoc rem concipiendi modo sponte elucet, quaecunque experimenta instituantur cum acu magnetica in \textit{unico} statu magnetico, ex his solis inclinationem derivari non posse, sed opus esse ut situs centri gravitatis \textit{veri} aliunde iam innotuerit. Hic situs stabiliri solet, antequam acus magnetismo imbuatur: sed hic modus parum tutus est, quum plerumque acus chalybea iam inter ipsam fabricationem magnetismum utut debilem assumat. Necessarium itaque est pro determinatione inclinationis, ut per mutationem idoneam status magnetici acus, \textit{alia} transpositio centri gravitatis eliciatur, quae quo a priori quam maxime diversa evadat, polos invertere oportebit, quo pacto transpositio duplex obtineri potest. Ceterum transpositio centri gravitatis vel in acubus dimensionum aptissimarum magnetismoque usque ad saturationem imbutis certum limitem transcendere nequit, qui (pro transpositione simplice) in nostris regionibus est circiter $0{,}4$ millimetri, et in regionibus, ubi vis verticalis maxima est, infra $0{,}6$ millimetri ma\SourceLeaf{96}net: unde simul intelligitur, quanta subtilitas mechanica in acubus ad inclinationem determinandam destinatis requiratur.

\Article{7}
Si corporis magnetici punctum aliquod $C$ fixum supponitur, ad aequilibrium requiritur et sufficit, ut planum per $C$, centrum gravitatis atque axem magneticum ductum cum plano meridiano magnetico coincidat, praetereaque momenta, quibus vis magnetica terrestris atque gravitas illud planum circa punctum $C$ vertere nituntur, se destruant: posterior conditio eo redit, ut denotante $T$ partem horizontalem vis magneticae terrestris, $i$ inclinationem axis magnetici ad planum horizontale, esse debeat $TM\sin i$ aequalis producto e pondere corporis in distantiam centri gravitatis transpositi $B'$ a recta verticali per $C$ ducta: manifesto haec distantia esse debet a parte australi vel boreali, prout $i$ est elevatio vel depressio, et pro $i=0$, $B'$ in ipsa ista recta verticali. Quodsi iam corpus circa hanc verticalem ita motum fuerit, ut axis magneticus pervenerit in planum verticale, cuius azimuthum magneticum, i. e. angulus cum parte boreali meridiani magnetici, (ad lubitum vel versus orientem vel occasum pro positivo acceptum) sit $=u$, magnetismus terrestris exseret vim ad corpus circa axem verticalem vertendum, i. e. ad angulum $u$ minuendum, cuius momentum erit $=TM\cos i\sin u$, corpusque circa hunc axem oscillationes faciet, quarum duratio per methodos notas calculari potest. Scilicet denotando per $K$ momentum inertiae corporis respectu axis oscillationis (i. e. aggregatum molecularum ponderabilium multiplicatarum per quadrata distantiarum ab axe), et pro more per $\pi$ semicircumferentiam circuli pro radio $=1$, erit tempus unius oscillationis infinite parvae $=\pi\BareRoot{\frac{K}{TM\cos i}}$. siquidem quantitatibus $T,M$ subest unitas virium acceleratricium ea, quae in unitate temporis gignit velocitatem $=1$: reductio oscillationum finitarum ad infinite parvas simili modo ut pro oscillationibus penduli calculari poterit. Quodsi igitur tempus unius oscillationis infinite parvae ex observationibus erutum est $=t$, habebimus $TM=\frac{\pi\pi K}{tt\cos i}$, adeoque, si quod semper abhinc subintelligimus, corpus ita suspensum est, ut axis magneticus sit horizontalis
\[
TM=\frac{\pi\pi K}{tt}
\]
Si magis placet, gravitatem pro unitate virium acceleratricium adoptare, illum
